\documentclass[a4paper, 12pt, french]{article}

% ========================
% IMPORTATION DES PACKAGES
% ========================

\usepackage[T1]{fontenc}
\usepackage[utf8]{inputenc}
\usepackage[margin=1.5cm]{geometry}                             \usepackage{lmodern}
\usepackage{theorem}                                            \usepackage{amsfonts, amsmath, amssymb}
\usepackage{babel}

\renewcommand{\familydefault}{\sfdefault}  

\theoremstyle{break}
\theoremheaderfont{\scshape}
\theorembodyfont{\upshape}

\newtheorem{exer}{Exercice}[section]
\newtheorem{sol}{Solution}[section]


\title{EXERCICE DE PROBABILITÉ}
\date{\today}
\author{}

\begin{document}

\maketitle

Dans cette document, j'apprend la probabilité en résoudrons des exercice et des preuve du cours de probabilité tout en développons mes competences en latex.\\
Chaque exercice et les énoncés dans cette document est démontrer avec mon propre raisonnement surtout la démonstration mais si cela peut vous aider, je vous invite à l'explorer.

\section{Espace probabilisé}
\begin{exer}
	 Montrer que si A et B sont indépendant, alors $A^{c}$ et B le sont aussi, ainsi que $A^{c}$ et $B^{c}$.
\end{exer}

\begin{sol}
	Soient A et B deux évenement et $\Omega$ l'ensemble de tous les évenements. Supposons que A et B sont indépendant. On sait que $\Omega = A \cup A^{c}$, on peut donc écrire B comme suit $B = B \cap \Omega = (B \cap A) \cup (B \cap A^{c})$.Comme $(B \cap A)$ et $(B \cap A^{c})$ sont disjoints alors d'après la $\sigma$-additivité, on a  $P(B) = P(B \cap A) + P(B \cap A^{c}$. Par hypothèse on a $P(A \cap B) = P(A)P(B)$. Ainsi, on a $P(B \cap A^{c} = P(B) - P(A)P(B)$ cela implique que $P(B \cap A^{c}) = P(B)(1-P(A))$ cela implique encore que  $P(B \cap A^{c}) = P(B)P(A^{c})$.\\ D'où B et $A^{c}$ sont indépendant.\\ \\
	Pour l'autre on utilisant la loi de Morgan, on a $A^{c} \cap B^{c} = (A \cup B)^{c}$. Ainsi $P(A^{c} \cap B^{c}) = 1 - P(A \cup B) = 1 - P(A) - P(B) + P(A)P(B) = 1-P(A) -P(B)(1 - P(A)) = P(A^{c}) -P(B)P(A^{c}) = P(A^{c})(1-P(B)) = P(A^{c})P(B^{v})$.\\ D'où $A^{c}$ et $B^{c}$ sont indépendant.
\end{sol} 

\begin{exer}
	Montrer que si P($\bigcap_{i \in \mathbb{N}} A_{i}$) > 0, alors $P(\bigcap_{i=1}^{n}) = P(A_{1}) P(A_{2}|A_{1}) P(A_{3}| A_{2} \cap A_{1}) \cdots P(A_{n}|\bigcap_{i=1}^{n-1} A_{i})$.
\end{exer}
\begin{sol}
	Raisonnons par récurrence, pour n = 2 on a :\\
	$P(A_{1} \cap A_{2}) = P(A_{1})P(A_{2}|A_{1})$, c'est vrai car c'est la formule de labprobabilité conditionnelle.\\
	Supposons que $P(\bigcap_{i=1}^{n}) = P(A_{1}) P(A_{2}|A_{1}) P(A_{3}| A_{2} \cap A_{1}) \cdots P(A_{n}|\bigcap_{i=1}^{n-1} A_{i})$ est vrai à l'ordre n et montrons qu'il est aussi vrai à l'ordre n+1.\\
	On a $P(\bigcap_{i=1}^{n+1} A_{i}) = P(\bigcap_{i=1}^{n}  A_{i} \cap A_{n+1})$. On a $P(A_{n+1}| \bigcap_{i=1}^{n} A_{i})=  \frac{P(\bigcap_{i=1}^{n+1} A_{i})}{P(\bigcap_{i=1}^{n} A_{i})}$. Ainsi $P(\bigcap{i=1}^{n+1} A_{i}) = P(A_{n+1}| \bigcap_{i=1}^{n} A_{i}) P(\bigcap_{i=1}^{n} A_{i})$. Or par hypothèse des récurrence, $P(\bigcap_{i=1}^{n} A_{i}) = P(A_{1}) P(A_{2} | A_{1}) \cdots P(A_{n}|\bigcap_{i=1}^{n-1} A_{i})$. Donc $P(\bigcap{i=1}^{n+1} A_{i})  =P(A_{2}|A_{1}) \cdots P(A_{n+1}| \bigcap_{i=1}^{n} A_{i}))$.\\
	On peut en conclure alors que $P(\bigcap_{i=1}^{n}) = P(A_{1}) P(A_{2}|A_{1}) P(A_{3}| A_{2} \cap A_{1}) \cdots P(A_{n}|\bigcap_{i=1}^{n-1} A_{i})$.

\end{sol}
\end{document}
